%% ma: L10-B14-0100
%% lop: 10
%% chuong: 5
%% bai: 14
%% dang: 10.14.8
%% loai: TLN
%% muc: VD
%% dap_an: "1,89"
%% nguon: "chuyen-de-cac-so-dac-trung-cua-mau-so-lieu-khong-ghep-nhom-toan-10-chuong-trinh-moi.pdf (GV. Phan Nhật Linh) (P05), Bài 03 Dạng 2 Phần III Câu 6"
%% kien_thuc: "Lớp 10 Bài 13 (số trung bình), Bài 14 (phương sai, độ lệch chuẩn)"
%% trang_thai: chinh-thuc
%% ghi_chu: "Bổ sung yêu cầu làm tròn đến hàng phần trăm. Tự duyệt theo yêu cầu giáo viên 22/10/2026."
\begin{ex}
Điều tra số sách tham khảo môn Toán của $30$ học sinh ở một lớp 10, ta thu được mẫu số liệu:
\begin{center}
\begin{tabular}{|c|c|c|c|c|c|c|c|c|c|}\hline
$6$ & $1$ & $4$ & $6$ & $7$ & $2$ & $7$ & $5$ & $3$ & $5$\\\hline
$5$ & $2$ & $2$ & $2$ & $2$ & $4$ & $3$ & $2$ & $1$ & $3$\\\hline
$7$ & $6$ & $6$ & $3$ & $3$ & $4$ & $7$ & $4$ & $3$ & $2$\\\hline
\end{tabular}
\end{center}
Xác định độ lệch chuẩn của mẫu số liệu (làm tròn đến hàng phần trăm).
\shortans{1,89}
\loigiai{Bảng tần số: giá trị $1$; $2$; $3$; $4$; $5$; $6$; $7$ có tần số $2$; $7$; $6$; $4$; $3$; $4$; $4$.

$\sum n_ix_i=117$; $\sum n_ix_i^2=563$.

$s^2=\dfrac{563}{30}-\left(\dfrac{117}{30}\right)^2\approx3{,}557$; $s\approx1{,}89$.}
\end{ex}
